\begin{table}[t]
\centering
\caption{Standardized Effect Sizes}
\label{tab:sde}
\begin{adjustbox}{max width=\textwidth}
\begin{tabular}{lcccccc}
\toprule
Outcome & $\hat{\beta}$ & SE & SD($Y$) & SDE & SE(SDE) & Classification \\
\midrule
\multicolumn{7}{l}{\textit{Panel A: Pooled}} \\
LBW rate (pooled TWFE) & 0.0510 & 0.0304 & 2.0318 & 0.0251 & 0.0150 & Small positive \\
LBW rate (CS-DiD) & 0.0289 & 0.0216 & 2.0318 & 0.0142 & 0.0106 & Small positive \\
\midrule
\multicolumn{7}{l}{\textit{Panel B: Heterogeneous}} \\
LBW rate (high capacity) & 0.0107 & 0.0553 & 2.0318 & 0.0053 & 0.0272 & Small positive \\
LBW rate (high poverty) & 0.1260 & 0.0753 & 2.0318 & 0.0620 & 0.0371 & Moderate positive \\
LBW rate (rural) & 0.0733 & 0.0548 & 2.0318 & 0.0361 & 0.0269 & Small positive \\
\bottomrule
\end{tabular}
\end{adjustbox}
\begin{tablenotes}
\item \textit{Notes:} \textbf{Country:} United States. \textbf{Research question:} Does coal plant pollution control mandated by the EPA Mercury and Air Toxics Standards reduce county-level low birth weight rates near affected facilities? \textbf{Policy mechanism:} MATS required approximately 296 coal-fired power plants to install mercury, acid gas, and fine particulate controls in three staggered compliance waves (April 2015, 2016, 2017), reducing mercury emissions by 86\% and acid gas HAPs by 96\% relative to pre-regulation levels. \textbf{Outcome definition:} Share of births below 2,500 grams (low birth weight rate in percentage points) from County Health Rankings annual data. \textbf{Treatment:} Binary; county within 50 miles of a MATS-affected coal plant, with timing based on the nearest plant's compliance wave. \textbf{Data:} County Health Rankings 2012--2020 (underlying birth data approximately 2008--2016), EIA Form 860 for plant compliance timing and coordinates, Census SAIPE for economic controls; 2,765 counties, 24,287 county-year observations. \textbf{Method:} Callaway--Sant'Anna (2021) staggered DiD with never-treated comparison group and doubly-robust estimation; TWFE with county and year fixed effects; standard errors clustered at county level. \textbf{Sample:} Continental U.S. counties; exposed group defined as counties within 50 miles of at least one coal plant; excluded Alaska, Hawaii, and territories. SDE $= \hat{\beta} / \text{SD}(Y)$ where SD($Y$) is the pre-treatment standard deviation of the low birth weight rate. Classification refers to magnitude, not statistical significance: Large ($|$SDE$| > 0.15$), Moderate ($0.05$--$0.15$), Small ($0.005$--$0.05$), Null ($< 0.005$).
\end{tablenotes}
\end{table}

